% chapters/02/section-05.tex
\section{Keputusan dan Navigasi pada NPC}
\label{sec:keputusan-navigasi}

Bagian ini membahas landasan teori bagi Blok~II, yaitu cara membuat objek yang
dikendalikan komputer---lazim disebut NPC, kependekan dari
\istilah{non-player character}---berperilaku masuk akal di dalam permainan.
Landasan ini dipakai pada Bab~VIII sampai Bab~XI.

\subsection{Dua Pertanyaan yang Berbeda}

Perilaku NPC yang tampak cerdas sebenarnya merupakan hasil kerja dua lapisan
yang menjawab dua pertanyaan berbeda.

\begin{itemize}
  \item \textbf{Lapisan keputusan} menjawab pertanyaan ``apa yang sebaiknya
        saya lakukan sekarang?'' Jawabannya berupa niat, misalnya berpatroli,
        mengejar pemain, atau mundur.
  \item \textbf{Lapisan navigasi} menjawab pertanyaan ``bagaimana cara menuju
        ke sana?'' Jawabannya berupa rangkaian langkah atau arah gerak.
\end{itemize}

Memisahkan keduanya membawa manfaat besar. Keputusan dapat diubah tanpa
menyentuh cara berpindah, dan sebaliknya, cara berpindah dapat diganti tanpa
mengubah aturan perilaku. Pemisahan itu juga memudahkan pengujian, sebab
setiap lapisan dapat diperiksa secara terpisah. Kedua lapisan tersebut beserta
gelang umpan baliknya diperlihatkan pada Gambar~\ref{fig:lapisan-npc}.

\begin{figure}[htbp]
\centering
% chapters/02/figures/lapisan-npc.tikz
% Dua lapisan pengendalian NPC: keputusan (mesin keadaan) dan navigasi
% (pencarian jalur), dengan gelang umpan balik dari dunia permainan.
\begin{tikzpicture}[
  lapis/.style={
    rectangle, rounded corners=2pt, draw=ubPrimary, thick, fill=ubAlt,
    align=center, text width=62mm, inner sep=4pt, minimum height=11mm,
    font=\footnotesize
  },
  lapisnav/.style={lapis, draw=ubAccent, fill=white},
  dunia/.style={lapis, draw=ubGreen, fill=ubGreenBg},
  panah/.style={-{Stealth[length=2.4mm]}, thick, draw=ubAccent},
  gelang/.style={-{Stealth[length=2.4mm]}, thick, draw=ubWarn, dashed},
  label/.style={font=\scriptsize\itshape, ubPrimary, align=center}
]

\node[lapis] (persepsi) at (0,0)
  {\textbf{Persepsi}\\[1pt]\scriptsize masukan pemain dan keadaan dunia};
\node[lapis] (keputusan) at (0,-1.9)
  {\textbf{Lapisan keputusan}\\[1pt]\scriptsize mesin keadaan: diam, berpatroli, mengejar\\[1pt]\scriptsize menghasilkan \itshape tujuan};
\node[lapisnav] (navigasi) at (0,-3.8)
  {\textbf{Lapisan navigasi}\\[1pt]\scriptsize pencarian jalur pada peta berhalangan\\[1pt]\scriptsize menghasilkan \itshape rangkaian langkah};
\node[lapis] (aksi) at (0,-5.7)
  {\textbf{Aksi}\\[1pt]\scriptsize gerak, animasi, dan suara};
\node[dunia] (dunia) at (0,-7.4)
  {\textbf{Dunia permainan}\\[1pt]\scriptsize keadaan yang diperbarui};

\draw[panah] (persepsi) -- (keputusan);
\draw[panah] (keputusan) -- node[right, label, xshift=2mm] {niat} (navigasi);
\draw[panah] (navigasi) -- node[right, label, xshift=2mm] {jalur} (aksi);
\draw[panah] (aksi) -- (dunia);

\draw[gelang] (dunia.west) -- ++(-2.8,0) |- (persepsi.west);
\node[label, anchor=west, text width=28mm] at (-9.4,-3.6)
  {umpan balik: setiap langkah diamati kembali};
\end{tikzpicture}

\caption{Lapisan keputusan dan navigasi pada pengendalian NPC}
\label{fig:lapisan-npc}
\end{figure}

\subsection{Agen Reaktif dan Agen Berbasis Perencanaan}

Pada lapisan keputusan, terdapat dua pendekatan yang berbeda tingkat
kerumitannya.

\textbf{Agen reaktif} menentukan tindakannya semata-mata berdasarkan keadaan
yang teramati saat itu. Contohnya, NPC yang bergerak lurus ke arah pemain
selama pemain berada dalam jarak tertentu \cite{reintech-game-ai}. Cara ini
mudah diterapkan dan murah dihitung, tetapi mudah dikelabui: begitu pemain
berlindung di balik dinding, NPC akan terus menabrak dinding itu.

\textbf{Agen berbasis perencanaan} memperhitungkan keadaan yang lebih luas,
misalnya mencari rute mengelilingi rintangan, sebelum memutuskan langkah
berikutnya. Hasilnya lebih meyakinkan, tetapi menuntut perhitungan tambahan
dan penataan yang lebih rumit.

Perlu ditegaskan bahwa pilihan di antara keduanya bukan soal benar atau salah,
melainkan soal kecukupan. Sebuah permainan dengan ruang terbuka dan tanpa
rintangan tidak memerlukan perencanaan; permainan berlatar labirin jelas
memerlukannya.

\subsection{Mesin Keadaan sebagai Model Keputusan}

Cara yang paling banyak dipakai untuk menyusun lapisan keputusan adalah mesin
keadaan (\istilah{finite state machine}). Gagasannya sederhana: NPC selalu
berada pada tepat satu keadaan, dan perpindahan antar-keadaan terjadi karena
adanya peristiwa atau syarat tertentu.

Tiga istilah yang perlu Anda kuasai sejak awal adalah keadaan, peristiwa, dan
transisi. \textbf{Keadaan} (\istilah{state}) menyatakan apa yang sedang
dilakukan NPC, misalnya berdiam, berpatroli, atau mengejar.
\textbf{Peristiwa} (\istilah{event}) menyatakan sesuatu yang terjadi, misalnya
pemain memasuki jarak pandang. \textbf{Transisi} adalah aturan yang menyatakan
perpindahan dari satu keadaan ke keadaan lain ketika suatu peristiwa terjadi.

Kekuatan pendekatan ini terletak pada batasan ``tepat satu keadaan pada satu
waktu''. Batasan itu menghapus kebutuhan memikirkan kombinasi perilaku yang
saling bertabrakan, misalnya NPC yang mengejar dan berpatroli sekaligus.
Susunan keadaan yang lazim pada permainan adalah diam, berpatroli, mengejar,
serta menyerang atau menghindar \cite{santos-fsm}.

Pendekatan ini juga dipakai di luar permainan; Unity menyediakan
perkakas bawaan untuk menyusunnya secara visual, yang dapat Anda lihat pada
Bab~XII \cite{unity-leveldesign}. Mengerjakannya sendiri terlebih dahulu pada
Bab~IX bertujuan agar Anda memahami aturan yang mendasarinya, bukan sekadar
memindahkan kotak pada layar.

\subsection{Pencarian Jalur sebagai Persoalan Graf}

Pada lapisan navigasi, persoalan mencari rute dinyatakan sebagai pencarian pada
graf. Peta dibagi menjadi petak-petak, setiap petak yang dapat dilalui menjadi
simpul (\istilah{node}), dan hubungan antara dua petak yang bersebelahan
menjadi sisi (\istilah{edge}). Sisi dapat memiliki biaya, sehingga jarak yang
ditempuh dapat dibedakan menurut tingkat kesulitan medannya.

Pencarian jalur bertujuan menemukan rangkaian sisi yang menghubungkan tempat
asal dengan tempat tujuan dengan biaya sekecil mungkin. Dua sifat yang perlu
dibedakan pada setiap kaidah pencarian adalah:

\begin{itemize}
  \item \textbf{kelengkapan}, yaitu jaminan bahwa jalan akan ditemukan apabila
        memang ada, dan
  \item \textbf{keoptimalan}, yaitu jaminan bahwa jalan yang ditemukan adalah
        yang termurah.
\end{itemize}

Kaidah pencarian yang mempertimbangkan perkiraan jarak ke tujuan, disebut
heuristik, dapat jauh lebih cepat daripada pencarian yang menelusuri semua
kemungkinan secara buta. Perkiraan itu harus bersifat tidak melebihkan jarak
sebenarnya agar keoptimalan hasil tetap terjaga \cite{munir-astar}. Analisis
dan penerapannya, termasuk pengaruh pilihan heuristik terhadap jumlah simpul
yang diperiksa, dibahas pada Bab~X \cite{santos-pathfinding}.

\subsection{Anggaran Waktu dan Biaya Perhitungan}

Satu hal yang membedakan Blok~II dari dua blok lainnya adalah sifat
perhitungannya yang bergantung pada keadaan. Gerak objek biasa memerlukan
waktu yang tetap, berapa pun keadaannya. Sebaliknya, pencarian jalur baru
dijalankan ketika tujuan berubah atau ketika jalur yang sedang ditempuh tidak
lagi berlaku, dan biayanya tumbuh seiring luas peta yang ditelusuri.

Karena itu, perancangan perilaku NPC tidak dapat dipisahkan dari anggaran
waktu satu bingkai yang telah dibahas pada Bagian~\ref{sec:model-waktu}.
Apabila pencarian jalur dijalankan setiap bingkai untuk setiap NPC, anggaran
tersebut akan habis dan laju bingkai menurun. Cara-cara menghematnya, misalnya
menjalankan pencarian secara berkala dan menyimpan hasilnya untuk dipakai
ulang, akan Anda susun pada Bab~XI. Contoh penerapan pada permainan yang lebih
besar tersedia pada \mbox{[R25]} dan \mbox{[R34]}.
